\section{\iflanguage{english}{Preliminaries}{Grundlagen}}\label{sec:preliminaries}

\subsection{Kissat}

Kissat ist ein hochmoderner SAT-Löser nach dem Prinzip des Conflict-Driven Clause Learning (CDCL), der 
von Armin Biere und seiner Forschungsgruppe an der Johannes Kepler Universität Linz entwickelt 
wurde. Seit seiner ersten Veröffentlichung im Jahr 2020 hat sich Kissat zu einem der erfolgreichsten 
Löser in den jüngsten SAT-Wettbewerben entwickelt, zahlreiche Medaillen gewonnen und bei einer Vielzahl 
von Benchmark-Familien eine außergewöhnliche Leistung gezeigt.

Der Name „Kissat“ leitet sich vom deutschen Wort „Kissat“ ab (eine spielerische Abwandlung von „kein SAT“)
und spiegelt die Philosophie des Solvers wider, auf eine saubere und effiziente Implementierung zu 
setzen. Kissat ist in der Programmiersprache C geschrieben und legt Wert auf Einfachheit, Leistung und 
Reproduzierbarkeit.

\subsubsection{Kissat Architektur}

Kissat implementiert den CDCL-Algorithmus mit mehreren ausgefeilten Optimierungen:

\begin{table}[htbp]
    \centering
    \begin{tabular}{lp{0.6\textwidth}}
        \toprule
        \textbf{Komponente} & \textbf{Beschreibung} \\
        \midrule
        Variablen Elimination & Entfernt Variablen durch auflösen der Klasuen, vereinfacht dadurch die Formel \\
        \midrule
        Ausschluss durch Subsumption & Entfernt Klauseln, die logisch aus anderen abgeleitet werden können \\
        \midrule
        Untersuchen & Durchsucht den Suchraum, um Ausbreitungswege zu finden \\
        \midrule
        Zwischenspeicherung & Speichert erlernte Informationen und nutzt sie nach einem Neustart wieder \\
        \midrule
        Mehrstufige Entscheidungsheuristik & Stellt zuletzt verwendete Variablen in den Vordergrund \\
        \bottomrule
    \end{tabular}
    \caption{Kissat Optimierungen}\label{tab:kissat_optimizations}
\end{table}

\newpage

\subsubsection{Kissat-Konfigurationsvarianten}

Kissat bietet zahlreiche Konfigurationsoptionen, mit denen sich sein Verhalten anpassen lässt. Die in 
dieser Arbeit verwendeten Varianten stellen unterschiedliche Kompromisse zwischen dem Aufwand für die 
Vorverarbeitung und der Rechengeschwindigkeit dar:

\begin{table}[htbp]
    \centering
    \begin{tabular}{lp{0.6\textwidth}}
        \toprule
        \textbf{Variante} & \textbf{Haupteigenschaften} \\
        \midrule
        kissat\_default & Standardkonfiguration mit ausgewogenen Einstellungen \\
        \midrule
        kissat\_noeliminate & Deaktiviert die Variableneliminierung (schnellere Vorverarbeitung, möglicherweise langsamere Lösung) \\
        \midrule
        kissat\_nostable & Deaktiviert den stabilen Modus (abweichende Strategie zur Variablenauswahl) \\
        \midrule
        kissat\_stable & Aktiviert den stabilen Modus für eine systematischere Suche \\
        \midrule
        kissat\_stable2 & Erweiterter Stabilitätsmodus mit zusätzlichen Heuristiken \\
        \midrule
        kissat\_target2 & Optimiert für eine bestimmte Problemklasse \\
        \bottomrule
    \end{tabular}
    \caption{Kissat-Konfigurationsvarianten}\label{tab:kissat_variants}
\end{table}

\subsubsection{Messung der Solver-Leistung}

Die Leistung eines Solvers wird in der Regel anhand der Gesamtlaufzeit oder der CPU-Zeit gemessen. In 
dieser Arbeit verwende ich die vom Feld „process-time“ von Kissat gemeldete Laufzeit, die die zur Lösung 
der Instanz aufgewendeten CPU-Sekunden angibt. Zeitüberschreitungen (z. B. 3600 Sekunden, 7200 Sekunden) 
deuten darauf hin, dass der Solver die Berechnung nicht innerhalb der zugewiesenen Zeit abgeschlossen hat.


\subsection{Satzilla}

Satzilla (Xu et al., 2008) ist ein wegweisendes System, das maschinelles Lernen auf das Problem der 
Solverauswahl anwendet. Die zentrale Erkenntnis besteht darin, dass SAT-Instanzen durch Merkmale 
charakterisiert werden können, die die Leistung des Solvers vorhersagen. Durch das Lernen aus 
historischen Laufzeitdaten kann Satzilla den besten Solver für bisher unbekannte Instanzen auswählen.\cite{xu2008satzilla}

\newpage

\subsubsection{Feature Extraction}

Satzilla extrahiert aus jeder CNF-Formel über 100 Merkmale, die in mehrere Kategorien unterteilt sind:

\begin{table}[htbp]
    \centering
    \begin{tabular}{lp{0.3\textwidth}p{0.4\textwidth}}
        \toprule
        \textbf{Kategorie} & \textbf{Features} & \textbf{Was sie messen} \\
        \midrule
        Basic & Anzahl der Variablen, Anzahl der Klauseln, Verhältnis & Größe und Dichte \\
        \midrule
        Variable-Clause Graph (VCG) & Mittelwert, Entropie, Min/Max & Verbindungsstruktur eines Graphen \\
        \midrule
        Clause Graph (CG) & Klauselverbindungen, Cluster & Wechselwirkungen zwischen Klauseln \\
        \midrule
        Variable Graph (VG) & Variable co-occurrence & Variable Beziehungen \\
        \midrule
        Horn Features & Horn-Klausel-Bruch, Eigenschaften der Horn-Variablen & Besondere Struktur \\
        \midrule
        Probing Features & Ausbreitungstiefe der Einheit, reduzierte Variablen & Eigenschaften des Suchraums \\
        \midrule
        Local Search & Leistungskennzahlen von GSAT/SAPS & Die Landschaft der lokalen Suche \\
        \midrule
        LP Features & Relaxationswerte der linearen Programmierung & Kontinuierliche Entspannung \\
        \bottomrule
    \end{tabular}
    \caption{SAT Instance Features Categories}\label{tab:sat_features}
\end{table}

Diese Merkmale werden von einem speziellen Merkmalsextraktor (Merkmalsexekutivdatei) berechnet, der die 
CNF-Datei analysiert und eine CSV-Datei ausgibt, die eine Zeile pro Instanz enthält.

\subsubsection{Machine Learning Model}

Satzilla nutzt „Random Forest“ (Breiman, 2001) als Lernalgorithmus. „Random Forest“ ist ein 
Ensemble-Verfahren, bei dem mehrere Entscheidungsbäume erstellt und deren Vorhersagen zusammengefasst 
werden. Wesentliche Vorteile für diese Anwendung:

\begin{itemize}
    \setlength{\itemsep}{0pt}
    \setlength{\parskip}{0pt}
    \setlength{\parsep}{0pt}
    \item Unterstützt gemischte Merkmalstypen (kontinuierlich, kategorial)
    \item Unempfindlich gegenüber irrelevanten Merkmalen (der Algorithmus ignoriert verrauschte Eingaben)
    \item Zeigt die Merkmalsbedeutung an (welche Merkmale am wichtigsten sind)
    \item Verhindert Überanpassung durch geeignete Hyperparameter
\end{itemize}

\newpage

Im ursprünglichen Satzilla sagt das Modell die Laufzeit direkt voraus. Für die Solverauswahl müssen wir 
jedoch nur wissen, welcher Solver am schnellsten ist – nicht die genaue Laufzeit. Diese Arbeit 
konzentriert sich auf die Genauigkeit der Solverauswahl.

\subsubsection{Training und Prediction Pipeline}

Die gesamte Satzilla-Pipeline besteht aus:

\[
\begin{tikzcd}
\text{CNF File} \arrow[r] & \text{Feature Extractor} \arrow[r] & \text{Features} \arrow[d] \\
& & \text{Trained Model} \arrow[d] \\
\text{Runtime Data (training)} \arrow[r] & \text{Training} \arrow[r] & \text{Predictions (which solver is fastest)}
\end{tikzcd}
\]


\textbf{1. Training Phase:} Gegeben sei eine Menge von Trainingsinstanzen mit bekannten Laufzeiten für jeden Solver:
\begin{itemize}
    \setlength{\itemsep}{0pt}
    \setlength{\parskip}{0pt}
    \setlength{\parsep}{0pt}
    \item Features aus jeder CNF-Datei extrahieren
    \item Ermitteln Sie für jede Instanz, welcher Solver am schnellsten war
    \item Trainieren Sie einen Random-Forest-Klassifikator, um den schnellsten Solver vorherzusagen
\end{itemize}

\textbf{2. Prediction Phase:} Für eine neue CNF-Datei:
\begin{itemize}
    \setlength{\itemsep}{0pt}
    \setlength{\parskip}{0pt}
    \setlength{\parsep}{0pt}
    \item Die gleichen Features extrahieren
    \item Wende das trainierte Modell an, um den schnellsten Solver vorherzusagen
    \item Führen Sie nur diesen Solver aus (das spart Zeit)
\end{itemize}

\newpage

\subsubsection{Warum Satzilla nach wie vor von Bedeutung ist}

Obwohl Satzilla bereits 2008 eingeführt wurde, sind seine Grundsätze nach wie vor aktuell:

\begin{itemize}
    \setlength{\itemsep}{0pt}
    \setlength{\parskip}{0pt}
    \setlength{\parsep}{0pt}
    \item Rechenintensivitätsarme Merkmale (Sekunden pro Instanz) im Vergleich zu rechenintensiven Solver-Läufen (Minuten bis Stunden)
    \item Einmalige Trainingskosten, die über viele Vorhersagen verteilt abgeschrieben werden
    \item Transparent und nachvollziehbar (die Merkmalsbedeutung zeigt, worauf es ankommt)
    \item Anpassbar an neue Solver-Sätze (wie in dieser Arbeit gezeigt)
\end{itemize}